\section{Création du Bloc QuadRAM}

\subsection{Conception et validation du décodeur d'adresses}

Un décodeur d'adresses générique a été conçu pour piloter les signaux \texttt{cs} des blocs RAM à partir des bits de poids fort de l'adresse. Son fonctionnement a été validé par simulation (Figure~\ref{fig:decoder}).

\begin{lstlisting}[caption={Décodeur d'adresses générique}, label={lst:decoder}]
entity decoder is
  generic(SEL_LENGHT : INTEGER := 2);
  port (
    address : in  unsigned((SEL_LENGHT - 1) downto 0);
    cs      : out unsigned(((2**SEL_LENGHT) - 1) downto 0)
  );
end entity;

architecture RTL of decoder is
begin
  process(address) is
  begin
    cs <= (others => '0');
    cs(to_integer(address)) <= '1';
  end process;
end architecture;
\end{lstlisting}

\begin{figure}[H]
    \centering
    \includegraphics[width=1\textwidth]{2_decoder_modified.png}
    \caption{Chronogramme du décodeur d'adresses}
    \label{fig:decoder}
\end{figure}

Ainsi, sur le chronogramme de la Figure~\ref{fig:decoder}, le vecteur de sortie \texttt{cs} évolue en fonction de la valeur du signal d’entrée \texttt{address}. Pour une adresse codée sur deux bits, quatre combinaisons sont possibles : \texttt{00}, \texttt{01}, \texttt{10} et \texttt{11}. Chacune active un bit de \texttt{cs}, tandis que les autres restent à zéro. Ce comportement permet de sélectionner le bloc mémoire correspondant à l’adresse fournie. La simulation confirme le bon fonctionnement du décodeur et son utilisation dans l’architecture de la \texttt{quadram}.


\subsection{Instanciation de 4 blocs SyncRam}

L'entité \texttt{QuadRAM} intègre un décodeur d'adresses et une boucle \texttt{for...generate} pour instancier les blocs mémoire. Le décodeur analyse les bits de poids fort de l'adresse pour piloter individuellement les signaux \texttt{cs} de chaque RAM, tandis que les autres signaux (horloge, écriture et données) sont mutualisés.
\\

\begin{lstlisting}[caption={QuadRAM}, label={lst:quadram}]
DECODER : entity work.decoder(RTL)
generic map (
    SEL_LENGHT => 2
)
port map (
    cs      => cs_array,
    address => address(9 downto 8)
);

GenRam : for i in 0 to (2**2) - 1 generate 
    Ram : entity work.sync_ram(RTL)
    generic map (
        DataLenght => 8,
        AddrLenght => 8
    )
    port map (
        cs      => cs_array(i),
        clock   => clock,
        we      => we,
        address => address(7 downto 0),
        datain  => datain,
        dataout => dataout
    );
end generate;
\end{lstlisting}

\subsection{Généralisation de l'architecture}

Pour généraliser l'architecture de la RAM, nous avons ajouté des généricités pour la taille des données et des adresses. Cela permet de définir la taille de la RAM en fonction des besoins spécifiques du projet.
\\

\begin{lstlisting}[caption={QuadRAM Généralisée}, label={lst:quadram_gene}]
	generic(ADDR_SELECT : INTEGER := 2;
		    ADDR_LENGHT : INTEGER := 10;
		    DATA_LENGHT : INTEGER := 8
	);
\end{lstlisting}

\subsection{Validation de la QuadRAM}

Afin de valider le fonctionnement de la QuadRAM, nous avons réalisé une simulation en utilisant un test bench. Le chronogramme obtenu (Figure~\ref{fig:quadram}) montre que la QuadRAM fonctionne correctement, avec les signaux \texttt{cs} activés pour chaque bloc RAM en fonction des bits de poids fort de l'adresse.
\\

\begin{figure}[H]
    \centering
    \includegraphics[width=1\textwidth]{2_QUADRAM_modified.png}
    \caption{Chronogramme de la QuadRAM}
    \label{fig:quadram}
\end{figure}

La Figure~\ref{fig:quadram} présente le chronogramme de validation de la \texttt{quadram}, composée de quatre blocs RAM pilotés par un décodeur d’adresses. Les 2 bits de poids fort de \texttt{address} activent le bit correspondant de \texttt{cs\_array}, tandis que les 8 bits de poids faible définissent l’adresse interne de la banque sélectionnée. Ainsi, une seule RAM est activée parmi les quatre instances. Cette simulation vérifie le décodage et la sélection des mémoires, nécessaires à l’utilisation de la \texttt{quadram} comme mémoire modulaire.